\chapter{Polynomial methods and arithmetic applications}\label{chap:M07.19}

% Symmetric polynomials
\KEntry{K07.19.symmetric_functions}{Elementary symmetric polynomials}
Let $R$ be a commutative unital ring. In $R[X_1,\ldots,X_n]$ put
\[
 e_j=\sum_{1\le i_1<\cdots<i_j\le n}X_{i_1}\cdots X_{i_j}
 \quad(1\le j\le n),\qquad e_0=1.
\]
Permuting variables gives ring automorphisms: it permutes exponent
vectors in the finite coefficient expansion, and respects their sums
under multiplication. A polynomial is symmetric if all these
automorphisms fix it. The fixed polynomials form a subring, since the
automorphisms preserve subtraction, multiplication and one.
Expansion by choosing either $T$ or $-X_i$ in each factor gives
\begin{equation}\label{eq:symmetric-vieta}
 \prod_{i=1}^n(T-X_i)=\sum_{j=0}^n(-1)^j e_jT^{n-j}.
\end{equation}
This is a formal polynomial identity over every $R$, not an assertion
about polynomial functions. The empty case has product $1$.

\KEntry{K07.19.symmetric_polynomial_theorem}{The fundamental symmetric-polynomial theorem}
\begin{theorem}\label{thm:symmetric-polynomials}
For every symmetric $f\in R[X_1,\ldots,X_n]$ there is a unique
$G\in R[T_1,\ldots,T_n]$ with $f=G(e_1,\ldots,e_n)$.
Equivalently the substitution homomorphism identifies the polynomial
ring on $n$ symbols with the symmetric subring. No characteristic
restriction or division by $n!$ is required.
\end{theorem}
\begin{proof}
The zero ring and $n=0$ give the assertion directly, so assume $R$
nontrivial and $n>0$. Order exponent tuples lexicographically, with
$X_1$ first. For nonzero symmetric $f$, let $cX_1^{a_1}\cdots X_n^{a_n}$
be its largest monomial. Then $a_1\ge\cdots\ge a_n$:
if $a_i<a_j$ with $i<j$, swapping those variables gives a larger
monomial with the same nonzero coefficient, a contradiction.
The leading monomial of $e_j$ is $X_1\cdots X_j$, with coefficient
one. A product of polynomials with such leading coefficients has
the product leading monomial: lexicographic order respects addition
of tuples, and the maximal term has no competing contribution.
Consequently
\[
 E=e_1^{a_1-a_2}e_2^{a_2-a_3}\cdots e_{n-1}^{a_{n-1}-a_n}e_n^{a_n}
\]
has leading monomial $X_1^{a_1}\cdots X_n^{a_n}$ with coefficient
one. Subtract $cE$. The result remains symmetric, its leading
monomial decreases, and no monomial of total degree above $\deg f$
has been introduced: $E$ is homogeneous of degree $\sum a_i$.
There are finitely many tuples with total degree at most that bound.
Repeating therefore terminates and expresses $f$ as a polynomial
in the $e_j$.

For uniqueness, a monomial $T_1^{b_1}\cdots T_n^{b_n}$ maps to a
polynomial with leading exponent tuple
\[
 (b_1+\cdots+b_n,\ b_2+\cdots+b_n,\ \ldots,\ b_n)
\]
and leading coefficient one. Distinct $b$ give distinct tuples,
since adjacent differences recover $b$. In a nonzero finite sum
$G$, choose the image-leading tuple largest among its nonzero
terms. All other images have smaller leading tuples, so cannot
cancel its nonzero coefficient. Thus $G(e_1,\ldots,e_n)\ne0$,
proving injectivity and uniqueness, even when $R$ has zero divisors.
\end{proof}
Giving $T_j$ weight $j$ makes substitution preserve total degree.
The argument can be applied to homogeneous components and shows that
the unique expression of a homogeneous symmetric polynomial is
weighted homogeneous. This theorem concerns invariants under variable
permutations. The symmetric powers and symmetric multilinear forms
of D05, including \texttt{K05.15.symmetric} and
\texttt{K05.15.polynomial\_polarization}, are different constructions.

\KEntry{K07.19.symmetric_root_expressions}{Symmetric expressions in roots}
If a monic $f\in F[T]$ splits in an extension with roots
$\alpha_1,\ldots,\alpha_n$ repeated with multiplicity, comparison
in \eqref{eq:symmetric-vieta} identifies their elementary symmetric
expressions with the signed coefficients of $f$.
Every symmetric polynomial over $F$ in these roots therefore belongs
to $F$: apply Theorem~\ref{thm:symmetric-polynomials} and substitute
the coefficients into $G$. Separability is unnecessary.
The theorem does not say that an arbitrary expression in the roots,
or an arbitrary rational expression with a vanishing denominator,
belongs to $F$. Newton identities may be reused from
\texttt{K05.08.newton\_identities}; dividing by their integer
coefficients still requires those scalars to be invertible.

% A fixed convention for resultants
\KEntry{K07.19.resultant_definition}{The Sylvester determinant}
For degree bounds $n,m\ge0$, and polynomials $\deg f\le n$,
$\deg g\le m$, define $\operatorname{Res}_{n,m}(f,g)$ as follows.
Form the square matrix with $m$ rows the descending coefficient
vectors of $X^{m-1}f,\ldots,f$, followed by $n$ rows those of
$X^{n-1}g,\ldots,g$, in the columns
$X^{n+m-1},\ldots,1$. Its determinant is the resultant.
An empty block has no rows and the empty determinant is $1$.
The ordinary resultant of two nonzero polynomials uses their
actual degrees. Bounds will be displayed when coefficients may
specialize or the derivative may lose degree.

The transpose of this matrix represents
\[
 F[X]_{<m}\oplus F[X]_{<n}\longrightarrow F[X]_{<n+m},
 \qquad(A,B)\longmapsto Af+Bg
\]
in descending monomial bases. This follows by expanding each input
coefficient. Swapping its two row blocks proves
$\operatorname{Res}_{m,n}(g,f)=(-1)^{nm}\operatorname{Res}_{n,m}(f,g)$.
If $g=c\ne0$ is constant and $f$ has degree $n$, the actual-degree
resultant is $c^n$; this follows from its diagonal matrix.

\KEntry{K07.19.resultant_common_root}{The common-root criterion}
\begin{theorem}\label{thm:resultant-common-root}
For nonzero polynomials $f,g$ of positive degrees $n,m$ over a field,
their resultant is zero if and only if they have a common
positive-degree divisor, equivalently a common root in an algebraic
closure. It is also equivalent to a relation $Af+Bg=0$ with
$\deg A<m$, $\deg B<n$, and $(A,B)\ne(0,0)$.
\end{theorem}
\begin{proof}
The displayed square linear map is singular exactly when its
determinant is zero, equivalently when it has such a nonzero kernel
vector; these are the D05 determinant and injectivity criteria.
If $h$ is a common nonconstant divisor, use
$A=g/h$, $B=-f/h$. Conversely if $\gcd(f,g)=1$, Bezout gives
$uf+vg=1$. From $Af=-Bg$ we obtain $f\mid Bg$, hence $f\mid B$:
multiply the Bezout identity by $B$. The degree bound forces $B=0$,
and then the domain property forces $A=0$, a contradiction.
Finally a common nonconstant divisor has a root in a splitting
field. A common root would contradict the evaluated Bezout
identity if the gcd were one. This proves the last equivalence.
\end{proof}
The root is in an algebraic closure; it need not be in $F$.
For example, two copies of $X^2+1$ over $\Q$ have zero resultant
and no common rational root.

\KEntry{K07.19.resultant_bezout}{A polynomial certificate}
\begin{lemma}\label{lem:commutative-ring-adjugate}
For a square matrix $A$ of positive size over any commutative unital
ring, the cofactor adjugate satisfies
$\operatorname{adj}(A)A=A\operatorname{adj}(A)=(\det A)I$.
\end{lemma}
\begin{proof}
Define the determinant by its signed permutation sum and the
cofactors by deleting a row and a column, with the corresponding sign;
the determinant of the zero-size minor is one.
Grouping permutation terms by the entry in one fixed column gives
the cofactor expansion along that column. The $(i,j)$ entry of
$\operatorname{adj}(A)A$ is
$\sum_k C_{ki}A_{kj}$, the cofactor expansion of the determinant
with column $i$ replaced by column $j$. For $i=j$ this is $\det A$.
For $i\ne j$ the matrix has two equal columns; pair each permutation
with the permutation obtained by interchanging those columns.
Their products are equal by commutativity and their signs are
opposite, so their terms sum to zero. The pairing has no fixed point,
including in characteristic two. Thus every off-diagonal entry is
zero. Row expansion proves the other identity in exactly the same
way: the $(i,j)$ entry is the determinant with row $j$ replaced by
row $i$, hence is $\det A$ if $i=j$ and zero otherwise.
This proves the identities using no division or field hypothesis.
\end{proof}
This is the general-ring extension of the field adjugate calculation
in \texttt{K05.05.adjugate}; its broader hypotheses have just been proved.
\begin{proposition}\label{prop:resultant-bezout}
Over any commutative coefficient ring, for the fixed bounds above
there are polynomials $U,V$ with
\[
 \operatorname{Res}_{n,m}(f,g)=Uf+Vg,\qquad
 \deg U<m,\quad \deg V<n,
\]
where a negative degree bound requires a zero polynomial.
If $n+m=0$, instead the determinant is $1$ and this certificate
is not asserted.
\end{proposition}
\begin{proof}
Let $S$ be the Sylvester matrix and
$z=(X^{n+m-1},\ldots,1)^t$. Then $Sz$ is the column of the row
polynomials used to define $S$. The adjugate identity
$\operatorname{adj}(S)Sz=(\det S)z$ holds over the commutative
polynomial ring by Lemma~\ref{lem:commutative-ring-adjugate}.
The last component gives the displayed
combination of $X^{m-1}f,\ldots,f,X^{n-1}g,\ldots,g$.
Collecting coefficients gives $U,V$ and their bounds.
\end{proof}
This certificate survives every coefficient homomorphism and
proves the necessary elimination implication even if degrees drop.

\KEntry{K07.19.resultant_product_discriminant}{Products over roots and discriminants}
\begin{theorem}\label{thm:resultant-product}
Let $f$ have degree $n\ge1$ and leading coefficient $a\ne0$ over
a field, and let $\deg g\le m$. In a splitting field write its
roots as $\alpha_1,\ldots,\alpha_n$ with multiplicity. Then
\[
 \operatorname{Res}_{n,m}(f,g)=a^m\prod_{i=1}^n g(\alpha_i).
\]
\end{theorem}
\begin{proof}
First $f$ is monic. In the target coefficient space of the
Sylvester map use the basis
$X^{m-1}f,\ldots,f,X^{n-1},\ldots,1$.
Its matrix in descending monomials is triangular with diagonal
ones, hence has determinant one. Division by $f$ shows that the
Sylvester map in this basis has blocks
$\begin{pmatrix}I_m&*\\0&\mu_g\end{pmatrix}$,
where $\mu_g$ is multiplication by $g$ on $F[X]/(f)$.
Thus the resultant is $\det\mu_g$.
Extend the scalar field to a splitting field; determinants are
preserved by mapping their finite coefficient formula.
If $f=\prod_\alpha(X-\alpha)^{e_\alpha}$, CRT identifies the
quotient with the product of the corresponding primary quotients.
In the basis $1,(X-\alpha),\ldots,(X-\alpha)^{e_\alpha-1}$
of one quotient, multiplication by $g$ is triangular with every
diagonal entry $g(\alpha)$: write
$g=g(\alpha)+(X-\alpha)q$. Taking block determinants gives the
claimed product, with all multiplicities.
For general $f=a f_0$, its $m$ rows are scaled by $a$, so its
determinant is $a^m$ times the monic resultant.
\end{proof}

For $f$ of degree $n\ge1$ define
\[
 \operatorname{Disc}(f)
 =a^{2n-2}\prod_{i<j}(\alpha_i-\alpha_j)^2.
\]
The product is symmetric and independent of the ordering. It
vanishes precisely when a root is repeated. For a monic polynomial
Theorem~\ref{thm:symmetric-polynomials} makes it a polynomial in
its coefficients. For general $f$ the precise resultant identity is
\begin{equation}\label{eq:discriminant-resultant}
 \operatorname{Disc}(f)=
 (-1)^{n(n-1)/2}a^{-1}\operatorname{Res}_{n,n-1}(f,f').
\end{equation}
Indeed $f'(\alpha_i)=a\prod_{j\ne i}(\alpha_i-\alpha_j)$
by differentiating the split product, even if a root repeats.
Multiplying pairs gives
$\prod_i f'(\alpha_i)=(-1)^{n(n-1)/2}a^n
\prod_{i<j}(\alpha_i-\alpha_j)^2$; apply the product theorem
with bound $n-1$.

The bound is essential in positive characteristic, where
$\deg f'$ can be smaller than $n-1$. Using its actual degree
without the corresponding leading-coefficient adjustment can
give a false nonmonic formula. For monic $f$ the extra factor is
one. Degree one gives discriminant one. For $n>1$ and $f'=0$,
the padded determinant and discriminant are both zero.
The discriminant is a polynomial with integer coefficients in
the coefficients of a degree-$n$ polynomial: in the universal
Sylvester determinant for $f,f'$, setting the leading coefficient
to zero makes the highest column zero. Hence that determinant
is divisible by the leading-coefficient variable in the integer
polynomial ring. Dividing it there and using
\eqref{eq:discriminant-resultant} gives the required universal
coefficient polynomial, valid on specialization.

\KEntry{K07.19.resultant_specialization}{What elimination does and does not show}
Fixed-bound resultants commute with coefficient homomorphisms,
because determinants are finite integer polynomials in the entries.
When $n+m>0$, if a specialization has a common root, the certificate in
Proposition~\ref{prop:resultant-bezout} forces the specialized
resultant to vanish. Conversely, over a specialized field it suffices
that at least one positive specified degree survives. For if
$n>0$ and $\deg f=n$, Theorem~\ref{thm:resultant-product} expresses the
resultant as a nonzero scalar times the product of the values of $g$
at the roots of $f$, repeated with multiplicity. A zero product over
the algebraic closure has a zero factor, which is a common root.
This also allows $g$ to become zero or constant. If instead $m>0$
and $\deg g=m$, interchange the two row blocks of the Sylvester matrix:
this changes its determinant by $(-1)^{nm}$ and reduces to the same
argument. If both bounds are zero, their fixed-bound resultant is
$1$, even when both specialized polynomials are zero; this empty-matrix
case is excluded from the necessary-elimination implication above.
Without a degree guard the converse can fail. In $\Q[Y][X]$ take
\[
 f=YX+1,\qquad g=YX+2,\qquad
 \operatorname{Res}_{1,1}(f,g)=
 \det\begin{pmatrix}Y&1\\Y&2\end{pmatrix}=Y.
\]
At $Y=0$ the resultant vanishes but the constants $1,2$ have no
common root. No claim that every projected algebraic point lifts,
or that a projection equals its closure, follows from this
necessary elimination equation. General elimination ideals and
Groebner algorithms remain explicit advanced undergraduate D22 topics.

% Additive normal forms for rational functions
\KEntry{K07.19.partial_fractions}{Normalized partial fractions}
\begin{theorem}\label{thm:normalized-partial-fractions}
Let $F$ be a field. Write a rational function with reduced monic
denominator $q=\prod_{i=1}^r p_i^{e_i}$, where the $p_i$ are
distinct monic irreducibles and $e_i\ge1$. There are unique
$h\in F[X]$ and polynomials $a_{ij}$ with $\deg a_{ij}<\deg p_i$
such that
\[
 \frac fq=h+\sum_{i=1}^r\sum_{j=1}^{e_i}\frac{a_{ij}}{p_i^j}.
\]
For a reduced denominator, $a_{i,e_i}\ne0$ for every $i$.
\end{theorem}
\begin{proof}
Division separates $f=hq+r_0$, $\deg r_0<\deg q$.
Set $q_i=p_i^{e_i}$ and $Q_i=q/q_i$. Since $Q_i$ is coprime
to $q_i$, Bezout makes it invertible modulo $q_i$.
Choose the unique representative $A_i$ of
$r_0Q_i^{-1}$ modulo $q_i$ with $\deg A_i<\deg q_i$.
The polynomial $r_0-\sum_i A_iQ_i$ vanishes modulo every $q_i$,
so CRT makes it divisible by $q$. Its degree is less than
$\deg q$, forcing it to be zero. Thus the proper part is
$\sum_i A_i/q_i$.

Successive division of $A_i$ by the monic $p_i$ gives the unique
expansion
$A_i=\sum_{k=0}^{e_i-1}b_{ik}p_i^k$ with $\deg b_{ik}<\deg p_i$.
The quotient degree decreases by $\deg p_i>0$ each time;
the original bound makes the quotient after $e_i$ steps zero.
Uniqueness follows by reducing modulo $p_i$, then subtracting
the first digit and cancelling $p_i$ inductively.
Put $a_{ij}=b_{i,e_i-j}$ to obtain the asserted expression.

For uniqueness of the grouped expression, subtract two expressions
and clear $q$. The sum $\sum_i(A_i-A_i')Q_i$ has degree less
than $\deg q$, whereas a nonzero $(h-h')q$ has degree at least
$\deg q$. Hence $h=h'$. Reducing the remaining sum modulo $q_i$
and cancelling the unit $Q_i$ gives $A_i=A_i'$ by their degree
bounds. Digit uniqueness gives every $a_{ij}=a_{ij}'$.
If $a_{i,e_i}=b_{i0}$ were zero, $A_i$ would be divisible by
$p_i$, and the congruence $A_iQ_i\equiv f\pmod {q_i}$
would give $p_i\mid f$, contradicting reduction.
\end{proof}
For a polynomial $q=1$, the sum is empty; zero has $h=0$.
Monicity and the numerator degree bounds remove unit and polynomial
ambiguities. If a denominator has already been factored into real
linear/quadratic factors or complex linear factors, the same theorem
specializes directly; existence of those factorization patterns uses
the separately proved fundamental theorem of algebra.

% Exact counting of real roots
\KEntry{K07.19.cauchy_root_bound}{A finite bound for roots}
For $f=a_nX^n+\cdots+a_0$, $a_n\ne0$, $n\ge1$, put
$M=\max_{i<n}|a_i/a_n|$. Every complex root satisfies
\begin{equation}\label{eq:cauchy-root-bound}
 |z|<1+M.
\end{equation}
If $M=0$, the only possible root is zero. Otherwise a hypothetical
root with $t=|z|\ge1+M>1$ would satisfy
\[
 t^n\le M\sum_{i=0}^{n-1}t^i
 =M\frac{t^n-1}{t-1}\le t^n-1,
\]
a contradiction. This is a bound on any root, not an appeal to
existence of complex roots or to FTA.

\KEntry{K07.19.sturm_chain}{Squarefree reduction and the Sturm chain}
For nonconstant $f\in\R[X]$, let $d$ be a monic gcd of $f,f'$ and
put $s=f/d$. Then $s$ has exactly the distinct real roots of $f$
and is squarefree, with $\gcd(s,s')=1$.
Here is a verification without complex factorization.
Factor $f=c\prod p_i^{e_i}$ in the UFD $\R[X]$.
Each irreducible $p_i$ has $p_i'\ne0$ and $\deg p_i'<\deg p_i$
in characteristic zero, so $p_i\nmid p_i'$.
Writing $f=p_i^{e_i}h$ with $p_i\nmid h$ gives
$f'=p_i^{e_i-1}(e_i p_i'h+p_i h')$.
The parenthesis is nonzero modulo $(p_i)$, since its quotient
is a field and $e_i\ne0$. Thus the gcd has precisely the
exponents $e_i-1$. This proves the assertions about $s$ and its roots.

For this squarefree $s$, form
$f_0=s,f_1=s'$ and
$f_{i+1}=-\operatorname{rem}(f_{i-1},f_i)$,
stopping before a zero remainder. Division strictly decreases
nonzero remainder degree, so terminates. Gcds are unchanged by
the division recurrence; hence the last member is a nonzero
constant because $\gcd(s,s')=1$.
Consecutive members have no common real zero: a common zero
propagates through
\[
 f_{i-1}=q_i f_i-f_{i+1}
\]
to the nonzero final constant, which is impossible.

\KEntry{K07.19.sturm_theorem}{Sturm's theorem with endpoint conventions}
\begin{theorem}\label{thm:sturm}
At a real $x$, delete zero entries from
$f_0(x),\ldots,f_\ell(x)$ and let $V(x)$ be the number of
changes of sign in the remaining list. If $a<b$ and
$s(a)s(b)\ne0$, the number of distinct roots of $f$ in $(a,b)$ is
$V(a)-V(b)$.
\end{theorem}
\begin{proof}
All chain polynomials are nonzero, so their real zero sets are
finite by the polynomial root bound. Away from their union all
signs are locally constant. On a subinterval containing none
of these zeros each sign is constant: a change of sign of a
continuous nonvanishing real polynomial would contradict IVT.

At a zero $c$ of an internal member $f_i$, its neighbors are
nonzero and satisfy $f_{i-1}(c)=-f_{i+1}(c)$.
Near $c$ their signs stay opposite. Whether $f_i$ is positive,
negative or deleted at $c$, the two edges across this triple
contribute exactly one variation. Thus this zero changes no
total variation. If several internal members vanish at $c$,
their positions are nonadjacent. The pairs of adjacent edges
just considered are disjoint, so the same conclusion holds
simultaneously for all of them.

At a root $c$ of $s$, squarefreeness gives $s'(c)\ne0$.
Write $s=(X-c)u$; differentiating gives $u(c)=s'(c)$.
Near $c$, $u$ and $f_1=s'$ therefore have the same fixed sign.
Immediately to the left, $f_0=s$ has the opposite sign to
$f_1$; immediately to the right their signs agree.
The first edge loses exactly one variation. Any internal
zeros at this same point contribute no change by the previous
paragraph, and do not use this first edge.
Consequently $V$ drops exactly one at every root of $s$ and
never elsewhere.

Partition $[a,b]$ at the finitely many chain zeros in it.
Because neither endpoint is a zero of $f_0$, deleting internal
zero entries there has the same value as its nearby interior
variation: the opposite-neighbor argument still applies.
Adding all changes gives exactly $V(a)-V(b)$, one for each
root in $(a,b)$. The roots of $s$ are the distinct roots of $f$.
\end{proof}
For a nonzero constant polynomial the root count is zero without
a chain. The zero polynomial is excluded. The displayed formula
does not count roots at the endpoints; if such endpoints are
wanted one must use one-sided limits and explicitly add the
chosen endpoint roots. At $+\infty$ a member's sign is its
leading-coefficient sign, and at $-\infty$ it is that sign
times $(-1)^{\deg f_i}$. Leading terms dominate for sufficiently
large absolute value, so these sign lists give the variation
limits and their difference counts all distinct real roots.
Alternatively choose endpoints outside the explicit bound
\eqref{eq:cauchy-root-bound}. No real-closed-field generalization
is assumed in this proof.

% Gaussian arithmetic
\KEntry{K07.19.gaussian_euclidean}{Euclidean division in the Gaussian integers}
In $\Z[i]\subset\C$, conjugation gives the multiplicative norm
$N(a+bi)=a^2+b^2$. It is a positive integer for nonzero elements.
A unit must have norm one; these elements are exactly
$1,-1,i,-i$, and their conjugates give their inverses.
For $z,w\in\Z[i]$, $w\ne0$, choose integers nearest to the real
and imaginary coordinates of $z/w$, and let $q$ be their
Gaussian combination. Then $r=z-qw$ satisfies
$N(r)=N(w)|z/w-q|^2\le N(w)/2<N(w)$.
Here $N(r)=0$ permits the zero remainder. This proves the
Euclidean property, so the earlier Euclidean--PID--UFD theorems
apply to this concrete ring. No unproved ideal-arithmetic
classification is needed.

\KEntry{K07.19.gaussian_primes_two_squares}{Gaussian primes and two squares}
\begin{theorem}\label{thm:prime-two-squares}
A positive integer prime $p$ is a sum of two integer squares
if and only if $p=2$ or $p\equiv1\pmod4$.
Up to units the Gaussian primes are $1+i$, rational primes
$q\equiv3\pmod4$, and Gaussian integers of norm a rational
prime $p\equiv1\pmod4$.
\end{theorem}
\begin{proof}
For an odd prime, a sum of two squares, each zero or one modulo
four, can only be one modulo four; $2=1^2+1^2$.
If $p\equiv1\pmod4$, inverse-pairing in the field $\Z/p\Z$
gives Wilson's congruence $(p-1)!\equiv-1$: the only
self-inverse nonzero residues are $1,-1$, by
$(x-1)(x+1)=0$, and all others pair to one.
Pairing instead $k$ with $p-k$ gives
\[
 (p-1)!\equiv(-1)^{(p-1)/2}
 \left(((p-1)/2)!\right)^2.
\]
Thus some integer $x$ has $x^2\equiv-1\pmod p$.
In $\Z[i]$, $p\mid(x+i)(x-i)$ but $p$ divides neither factor,
as their imaginary coefficients are $1,-1$. Therefore $p$
is not Gaussian prime. In this UFD it is not irreducible,
so $p=uv$ with both factors nonunits. Their positive norms
multiply to $p^2$ and both exceed one, forcing both to equal
$p$. This supplies $p=a^2+b^2$.

Every Gaussian integer of rational-prime norm is irreducible:
in a proper factorization the norms would be integers greater
than one multiplying to that prime. The element $1+i$ has
norm two, and $2=-i(1+i)^2$.
For a rational prime $q\equiv3\pmod4$, a Gaussian proper
factorization would force a factor of norm $q$ as above,
contradicting the modulo-four obstruction. Thus $q$ is a
Gaussian prime.

Conversely let $\pi$ be a Gaussian irreducible, hence prime.
It divides the positive integer $N(\pi)=\pi\bar\pi$, and hence
some rational prime $p$ in its integer factorization.
If $p=\pi\beta$, then $N(\pi)N(\beta)=p^2$.
The nonunit norm $N(\pi)$ is $p$ or $p^2$.
In the first case $\pi$ has prime norm, and modulo four
forces $p=2$ or $p\equiv1\pmod4$; norm two elements are
associates of $1+i$. In the second case $\beta$ is a unit,
so $\pi$ is associate to $p$. The cases $p=2$ and
$p\equiv1\pmod4$ were just shown reducible, leaving precisely
$p\equiv3\pmod4$. This proves exhaustion.
\end{proof}
This is a prime representation theorem. It does not silently
assert the further all-integer two-square classification or
unique ideal factorization in arbitrary quadratic rings.

% A corrected polynomial Fermat application
\KEntry{K07.19.mason_stothers}{Polynomial abc in characteristic zero}
For nonzero $u\in F[X]$, let $\operatorname{rad}(u)$ be the
product of its distinct monic irreducible factors; a nonzero
constant has radical one. The following argument is supplied
independently from polynomial UFD and differentiation; it is
not attributed to the printed Aluffi Fermat assertion.
\begin{theorem}\label{thm:mason-stothers}
Over a characteristic-zero field, if nonzero pairwise-coprime
$a,b,c$ satisfy $a+b=c$ and are not all constant, then
\[
 \max(\deg a,\deg b,\deg c)\le
 \deg\operatorname{rad}(abc)-1.
\]
\end{theorem}
\begin{proof}
Put $W=a'b-ab'$. It is nonzero. Otherwise coprimality of $a,b$
would give $a\mid a'$; a nonconstant $a$ cannot divide its
nonzero smaller-degree derivative. Thus $a$ would be constant;
then $W=0$ gives $b'=0$, and characteristic zero makes $b$
and $c$ constant too, a contradiction.
If $p^e$ exactly divides $a$, differentiation shows
$p^{e-1}\mid a'$, hence $p^{e-1}\mid W$.
The same holds for a factor of $b$.
Since $c=a+b$, we also have $W=a'c-ac'$, giving the same
divisibility for a factor of $c$. Pairwise coprimality and
unique factorization therefore give
$abc/\operatorname{rad}(abc)\mid W$, up to a unit.
Taking degrees yields
\[
 \deg a+\deg b+\deg c-\deg\operatorname{rad}(abc)
 \le\deg W\le\deg a+\deg b-1.
\]
This bounds $\deg c$. Apply the identical argument to
$c+(-b)=a$ and to $c+(-a)=b$; nonzero constants and units
do not affect coprimality or radical degree. All three bounds
follow.
\end{proof}

\KEntry{K07.19.primitive_fermat}{Primitive polynomial Fermat}
\begin{corollary}\label{cor:primitive-polynomial-fermat}
Let $F$ have characteristic zero and $n\ge3$.
For nonzero pairwise-coprime $f,g,h\in F[X]$,
$f^n+g^n=h^n$ forces all three to be constant.
\end{corollary}
\begin{proof}
If $D=\max(\deg f,\deg g,\deg h)>0$, apply the theorem to
$f^n,g^n,h^n$. Their radical has degree at most
$\deg f+\deg g+\deg h\le3D$, so $nD\le3D-1$, impossible
for $n\ge3$. Thus $D=0$.
\end{proof}
Without coprimality the unqualified statement is false:
over $\C$, $f=g=X$, $h=2^{1/n}X$ is a nonconstant solution.
Clearing their common factor turns the triple into constants,
so does not preserve the premise that they are not all constant.
In general remove a common gcd from a nonzero solution.
The resulting three polynomials are pairwise coprime: a prime
dividing two would divide the third by the equation, contrary
to their common gcd being one. The corollary then says the
original triple is a common polynomial multiple of a constant
solution. This is the corrected conclusion corresponding to
Aluffi's polynomial application. Characteristic zero is material:
Frobenius powers produce different behavior in characteristic $p$.
